% =====================================================================
%  RESUM (català) · una pàgina · mateix contingut que el Resumen
% =====================================================================
\begin{otherlanguage}{catalan}
\chapter*{\resumenname}
\pdfbookmark[0]{\resumenname}{resum}

% Exemple d'estructura (substitueix-lo pel teu resum):
Reservar un espai d'estudi al campus obliga avui a consultar diverses fonts i
no permet saber si una sala és lliure en aquell moment. Aquest treball presenta
\emph{SalaLibre}, una aplicació mòbil i un servei web que mostren l'ocupació de
les sales en temps real i permeten reservar-les en dos passos.

La solució s'ha desenvolupat amb una metodologia iterativa en tres increments.
L'aplicació Android es comunica amb una API REST que desa les reserves en una
base de dades relacional i notifica els canvis als clients. S'han definit dotze
requisits funcionals i cinc de no funcionals, i cadascun s'ha validat amb, com
a mínim, un cas de prova.

Les proves de càrrega mostren que, amb memòria cau, la latència mitjana es manté
per sota dels \qty{100}{ms} fins a 400 usuaris concurrents, mentre que la versió
inicial superava aquest llindar a partir de 200. Una prova d'usabilitat amb vuit
estudiants va obtenir una puntuació SUS mitjana de 81 sobre 100. Les dades d'aquest exemple són fictícies.

\vspace{1em}
\noindent\textbf{\clavesname:} \tfgclavesCA
\end{otherlanguage}
\cleardoublepage
